\section{Introdução}
\label{sec:introducao}

Problemas de conectividade aparecem em diferentes contextos da computação sempre que é necessário manter grupos de elementos e responder, com rapidez, se dois deles pertencem ao mesmo conjunto. No caso de grafos, essa situação corresponde à conectividade dinâmica incremental: novas arestas são inseridas e deseja-se saber se dois vértices passaram a fazer parte da mesma componente conexa.

A estrutura de dados \emph{Disjoint Set Union} (DSU), também chamada de \emph{Union--Find}, oferece as operações \texttt{Find} (obtenção do representante de um conjunto) e \texttt{Union} (fusão entre conjuntos distintos). Quando implementada de forma ingênua, essa estrutura tende a perder desempenho à medida que o número de operações cresce. Já heurísticas clássicas, como união por rank e compressão de caminho, reduzem substancialmente esse custo e levam à análise amortizada associada a Tarjan, expressa em termos de $\mathcal{O}(\alpha(n))$, em que $\alpha$ é a inversa da função de Ackermann \cite{tarjan1975efficiency}.

Neste trabalho, comparam-se três variantes de DSU: uma implementação ingênua, uma com união por rank e outra com união por rank combinada à compressão de caminho. Para essa comparação, o código foi organizado de forma modular, instrumentado para registrar leituras e escritas nos vetores \texttt{parent} e \texttt{rank}, e executado sobre sequências reproduzíveis de operações. A análise de Tarjan \cite{tarjan1975efficiency} sustenta a interpretação teórica das heurísticas estudadas, enquanto os tempos, acessos e gráficos apresentados ao longo do texto vêm do benchmark descrito na Seção~\ref{sec:metodologia}. Como referência de aplicação, menciona-se o algoritmo de Kruskal, que utiliza a mesma interface de \emph{Union--Find} para detectar ciclos durante a construção de árvores geradoras mínimas.

As seções seguintes apresentam a fundamentação teórica, o posicionamento do trabalho, a metodologia adotada, a discussão dos resultados e, por fim, as conclusões do estudo.
